\begin{table}[!ht]
  \centering
  \caption{Summary of the comparison with related work.}
  \label{tab:revisoes-sintese}
  \footnotesize
  \setlength{\tabcolsep}{3pt}
  \renewcommand{\arraystretch}{1.15}
  \begin{tabularx}{\columnwidth}{@{}
    >{\raggedright\arraybackslash}p{2.1cm}
    >{\raggedright\arraybackslash}X
    >{\raggedright\arraybackslash}X
    @{}}
    \toprule
    Dimension & Existing reviews & Present review \\
    \midrule

    Scope
      & Examine robustness in FL and defenses against adversarial clients. RL- or bandit-based defenses are at most cataloged or identified as a future research direction.
      & Takes RL- and bandit-based defenses against adversarial clients in FL as its primary object of analysis (77 studies). \\

    Central question
      & Which defenses exist, how they perform under different attacks, and whether they remain effective under realistic conditions.
      & Whether performance attributed to the defense stems from learning rather than from the signal, decision granularity, or memory ($\Delta_{\mathrm{L}}$). \\

    Synthesis method
      & Compare complete mechanisms and show that simple baselines can compete with sophisticated defenses, without classifying which components each baseline holds constant.
      & Classifies each comparison by the components controlled by the baseline and distinguishes the presence of a comparison from the strength of the evidence it provides. \\

    Experimental quality
      & Some audits examine datasets, attacks, heterogeneity, and the completeness of experimental evaluations.
      & Jointly examines random seeds, statistical inference, the use of test data for reward computation or model selection, and the learning budget. \\

    \bottomrule
  \end{tabularx}
\end{table}
